% ---------------------------------------------------------------------
%  第 16 章 性能评估
% ---------------------------------------------------------------------
\chapter{性能评估}
\label{ch:performance}

本章报告 2026 年 11 月 19 日的端到端对比与内核扫描。
对照实现为 NESSi。两组原始记录位于 \code{.nessi/bench/}，
分别为当日 09:53 的端到端记录和 09:50 的内核扫描记录。

\section{测量方法}
\label{sec:bench-method}

测量遵循三条原则：
\begin{enumerate}[label=(\arabic*)]
    \item \textbf{交错}：被测实现与对照实现交替运行，避免机器负载
          漂移偏向某一方；
    \item \textbf{中位数}：每个用例取多次运行的中位数，而不是平均值，
          以降低异常值的权重；
    \item \textbf{同窗口}：所有运行在同一时间段内完成，并记录负载。
\end{enumerate}
机器为 Intel Core Ultra 5 125H；Rust 侧用 release 构建
（\code{opt-level=3}、ThinLTO、单代码生成单元），对照实现用
\code{-O3} 编译。每个用例运行 3 到 5 次。

\begin{lstlisting}[numbers=none]
cargo build --release --examples
python3 .nessi/bench/compare_e2e.py       # end-to-end suite
python3 .nessi/bench/compare_kernels.py   # kernel sweep
\end{lstlisting}

\section{端到端结果}
\label{sec:bench-e2e}

端到端套件包含 10 个算例，覆盖 Hubbard、Fal'kov--Kimball 与
Holstein 三个模型族。最新一轮（2026-11-19）的结果如下：

\begin{table}[htbp]
    \centering
    \caption[端到端耗时，单位秒，取中位数]{端到端耗时，单位秒，取中位数}
    \label{tab:e2e}
    \small
    \begin{tabular}{@{}lrrr@{}}
        \toprule
        算例 & 本库 & 对照 & 加速比 \\
        \midrule
        \code{hubbard_2b}            & 0.0383 & 0.1514 & 3.95$\times$ \\
        \code{hubbard_gw}            & 0.0993 & 0.2166 & 2.18$\times$ \\
        \code{hubbard_tpp}           & 0.0985 & 0.1914 & 1.94$\times$ \\
        \code{fkm_bethe_quench}     & 0.0487 & 0.1366 & 2.81$\times$ \\
        \code{holstein_bethe_migdal}& 0.1246 & 0.3589 & 2.88$\times$ \\
        \code{holstein_bethe_umig}  & 0.0492 & 0.2097 & 4.26$\times$ \\
        \code{holstein_nambu_migdal}& 0.3310 & 0.9096 & 2.75$\times$ \\
        \code{holstein_nambu_umig}  & 0.2921 & 0.7861 & 2.69$\times$ \\
        \code{holstein_impurity_migdal} & 0.1033 & 0.3258 & 3.15$\times$ \\
        \code{holstein_impurity_umig}   & 0.0472 & 0.1511 & 3.20$\times$ \\
        \midrule
        整轮中位数 & 1.3605 & 3.4242 & 2.52$\times$ \\
        \bottomrule
    \end{tabular}
\end{table}

逐项加速比用对照耗时中位数除以本库耗时中位数。整轮结果先将一轮内的各算例耗时相加，
再分别对两套实现取中位数，因此整轮中位数一般不等于逐项中位数之和。

\begin{figure}[htbp]
    \centering
    \begin{tikzpicture}
        \begin{axis}[
            width=0.92\textwidth,
            height=6.4cm,
            ybar,
            bar width=9pt,
            ymin=0, ymax=4.6,
            ylabel={加速比},
            symbolic x coords={2b,gw,tpp,fkm,hbm,hbu,hnm,hnu,him,hiu},
            xtick=data,
            xticklabel style={rotate=30, anchor=east, font=\scriptsize},
            nodes near coords,
            nodes near coords style={font=\scriptsize},
            enlarge x limits=0.08,
            grid=major,
            grid style={gray!25}
        ]
            \addplot[fill=kbblue!55, draw=kbblue] coordinates {
                (2b,3.95) (gw,2.18) (tpp,1.94) (fkm,2.81) (hbm,2.88)
                (hbu,4.26) (hnm,2.75) (hnu,2.69) (him,3.15) (hiu,3.20)
            };
        \end{axis}
    \end{tikzpicture}
    \caption[端到端加速比]{端到端加速比。柱子按表中算例顺序排列，逐项为 1.94 至 4.26 倍。整轮加速比为 2.52 倍。}
    \label{fig:e2e-speedup}
\end{figure}

\section{内核扫描}
\label{sec:bench-kernels}

端到端结果混合了模型组装、观测量与输出，不能直接反映卷积与求解器
的效率。内核扫描单独测量三个核心家族在 $d=1,2,4,8$ 上的耗时，
固定 \code{ntau=200}、\code{k=5}、\code{rounds=3}。比值定义为
本库耗时除以对照耗时，大于 $1$ 表示本库更慢。

\begin{table}[htbp]
    \centering
    \caption[内核耗时对比，比值为本库除以对照]{内核耗时对比，比值为本库除以对照}
    \label{tab:kernel-sweep}
    \small
    \begin{tabular}{@{}lrrrr@{}}
        \toprule
        内核 & $d=1$ & $d=2$ & $d=4$ & $d=8$ \\
        \midrule
        卷积 & 1.60 & 1.12 & 1.56 & -- \\
        Dyson & 1.27 & 1.20 & 1.38 & 0.83 \\
        VIE2 & 1.24 & 1.63 & 1.87 & -- \\
        \bottomrule
    \end{tabular}
\end{table}

\begin{figure}[htbp]
    \centering
    \begin{tikzpicture}
        \begin{axis}[
            width=0.9\textwidth,
            height=6.2cm,
            ybar,
            bar width=8pt,
            ymin=0, ymax=2.2,
            ylabel={耗时比},
             symbolic x coords={d1,d2,d4,d8},
             xtick={d1,d2,d4,d8},
             xticklabels={$1$,$2$,$4$,$8$},
             xlabel={块尺寸 $d$},
            enlarge x limits=0.15,
            legend style={font=\scriptsize, at={(0.5,1.02)},
                          anchor=south, legend columns=3},
            grid=major,
            grid style={gray!25}
        ]
            \addplot[fill=kbblue!55, draw=kbblue] coordinates {
                (d1,1.60) (d2,1.12) (d4,1.56)
            };
            \addplot[fill=kbgreen!45, draw=kbgreen] coordinates {
                (d1,1.27) (d2,1.20) (d4,1.38) (d8,0.83)
            };
            \addplot[fill=kborange!45, draw=kborange] coordinates {
                (d1,1.24) (d2,1.63) (d4,1.87)
            };
            \legend{卷积, Dyson, VIE2}
        \end{axis}
    \end{tikzpicture}
    \caption[内核耗时比]{以本库耗时除以对照耗时。数值大于 1 表示本库更慢。
             $d=8$ 组仅有 Dyson 测量。}
    \label{fig:kernel-ratio}
\end{figure}

本轮端到端加速比为 $2.52\times$。内核扫描的比值方向相反。卷积与 VIE2 在各自已测块尺寸上的耗时比均大于 $1$，Dyson 在 $d=8$ 时为 $0.83$。两组结果分别反映完整算例和单项内核的运行时间。
